\section{Case-Study-Dossier: Projekt SafeLend (Generiert mit Gemini)}

\subsection{Unternehmenskontext und Ausgangslage}
Die \enquote{SafeLend Bank} ist ein etabliertes Kreditinstitut,
das auf Konsumkredite für Privatkunden spezialisiert ist.
Die Kreditprüfung erfolgt bisher über ein hybrides Verfahren
aus statistischen Schwellenwerten und manueller Endprüfung
durch Kreditanalysten.
Durch den Wettbewerbsdruck digitaler Anbieter ist die mittlere
Bearbeitungszeit von 48 Stunden nicht mehr wettbewerbsfähig.
Projektziel ist eine Senkung der Durchlaufzeit auf unter 10 Minuten.

\subsection{Zielsetzung des Projekts}
Ziel ist die Entwicklung eines Machine-Learning-basierten
Scoring-Modells für automatisierte Kreditentscheidungen.
Das Modell soll auf Basis historischer Transaktionsdaten,
Auskunftei-Informationen und demografischer Merkmale
die Ausfallwahrscheinlichkeit
\enquote{Probability of Default} robuster prognostizieren
als das bisherige Verfahren.
Zusätzlich sollen operative Kosten reduziert
und Entscheidungsqualität sowie Skalierbarkeit erhöht werden.

\subsection{Datenbasis und technische Infrastruktur}
Als Datengrundlage dienen rund 500.000 historische Kreditfälle
aus den letzten zehn Jahren.

Interne Daten umfassen Kontobewegungen der letzten 12 Monate,
Gehaltszuflüsse, bisheriges Zahlungsverhalten
sowie Mahn- und Rückstandsverläufe.

Externe Daten umfassen Bonitätsinformationen von Auskunfteien
sowie makroökonomische Indikatoren
wie regionale Arbeitslosenquote und Zinsniveau.

Die technische Umsetzung erfolgt cloudbasiert.
Das Modell wird als API-Service in die bestehende Banking-App integriert,
um Echtzeitentscheidungen im Antragsprozess zu unterstützen.

\subsection{Rahmenbedingungen und Herausforderungen}
Das Projekt bewegt sich in einem volatilen und regulierten Umfeld.

Regulatorisch unterliegt die Bank den relevanten Aufsichtsanforderungen,
insbesondere den Anforderungen an Nachvollziehbarkeit,
Dokumentation und Kontrollierbarkeit von Kreditentscheidungen.

Marktseitig verändert sich das Kundenverhalten
unter volatilen Rahmenbedingungen.
Hieraus entsteht ein zentrales Risiko in Form von Data Drift,
wenn historische Trainingsdaten künftige Muster
nicht mehr ausreichend repräsentieren.

Zusätzlich ist sicherzustellen,
dass systematische Benachteiligungen
zum Beispiel über indirekte Merkmalszusammenhänge
früh erkannt und verhindert werden.

\subsection{Projektorganisation}
Das Vorhaben wird interdisziplinär umgesetzt.
Beteiligt sind Fachabteilung, IT, externe Data-Science-Rollen
sowie Risiko- und Compliance-Funktionen.
Die Gesamtleitung liegt beim Head of Digital Banking.
Die operative Abstimmung erfolgt in wöchentlichen Sprints.

Im Projektverlauf bestehen erkennbare Spannungen:
Kreditanalysten adressieren Auswirkungen auf bestehende Rollenbilder,
während Data-Science-Rollen zunächst stärker
auf Modellmetriken fokussieren.
Damit ist eine explizite Governance- und Kommunikationsstruktur
projektkritisch.

\subsection{Zeitplan und Ressourcen}
Der geplante Umsetzungszeitraum beträgt neun Monate,
unterteilt in eine Pilotphase
\enquote{Proof of Concept}
und einen stufenweisen Rollout.
Das Gesamtbudget beträgt 550.000 Euro
einschließlich externer Beratung und Cloud-Ressourcen.
Der Go-Live ist für das folgende Geschäftsjahr vorgesehen.

\subsection{Risikeregister mit Status-Historie}
Die folgende Übersicht dokumentiert die Carry-Forward-Logik des SafeLend-Projekts.
Die Spalte Status-Historie zeigt, wie sich einzelne Risiken über die DASC-PM-Phasen entwickeln.

\begin{longtable}{@{}p{3.8cm}p{3.4cm}p{2.0cm}p{2.1cm}p{3.8cm}p{3.9cm}@{}}
\caption[SafeLend-Risikoregister]{SafeLend-Risikoregister mit Status-Historie}
\label{tab:safelend_status_history}
\\
\toprule
\textbf{Risiko} & \textbf{KRI / Trigger} & \textbf{Eintritt} & \textbf{Aus\allowbreak wirkung} & \textbf{Status-Historie} & \textbf{Maßnahmen / Owner} \\
\midrule
\endfirsthead
\caption[]{SafeLend-Risikoregister mit Status-Historie (Fortsetzung)}\\
\toprule
\textbf{Risiko} & \textbf{KRI / Trigger} & \textbf{Eintritt} & \textbf{Aus\allowbreak wirkung} & \textbf{Status-Historie} & \textbf{Maßnahmen / Owner} \\
\midrule
\endhead
\midrule
\multicolumn{6}{r}{Fortsetzung auf der nächsten Seite}\\
\endfoot

\bottomrule
\endlastfoot

Zielschieflage und Governance-Lücken
& Dokumentationsvollständigkeit der Governance-Artefakte = 100\%
& 5
& 6
& Phase 1: Gelb \textrightarrow{} Phase 2: Gelb \textrightarrow{} Phase 3: Grün
& Verbindliche Freigaben, Risk Owner, Eskalationspfad / Projektleitung \\

Datenqualität und Bias in historischen Daten
& Pflichtfeldvollständigkeit $\geq$ 98\%, Ablehnungsquoten-Differenz $\leq$ 5\%
& 6
& 10
& Phase 2: Rot \textrightarrow{} Phase 3: Grün
& Data-Quality-Gates, Great Expectations, Fairness-Prüfungen / Data Science \& Compliance \\

Instabile Modellgüte über Segmente
& Segment-AUC $\geq$ 0.75, Performance-Gap $\leq$ 0.05
& 5
& 6
& Phase 3: Gelb \textrightarrow{} Phase 4: Gelb
& Robustheitsanalysen, SHAP/LIME, manuelle Fachprüfung / Data Science \\

Data Drift im Betrieb
& PSI $> 0.25$ oder Incident-Reaktionszeit $> 30$ Min.
& 7
& 8
& Phase 4: Gelb \textrightarrow{} Phase 5: Rot
& Monitoring-Takt anpassen, Re-Training, Fallback / IT \& Risk Management \\

\bottomrule
\end{longtable}
